\documentstyle[%


righttag%


,11pt%


,amssymb%


,russian%


,izvru%


]{amsart}





\newcommand{\tl}[3]{\medskip\noindent{\bf#1}\ #2 \dotfill #3 \par} 


\pagestyle{empty}


% \pagestyle{plain}


\begin{document}


\par


\vskip 2cm


\par


\bigskip





Данный номер является тематическим, посвященным разработке


теории и методов решения задач оптимизации


\par





\bigskip


\bigskip


\par


\centerline{\Large СОДЕРЖАНИЕ}


\bigskip


\bigskip


%%%%%%%%%% Use \newline %%%%%%%%%%%





\tl{Аксенюшкина Е.В.}


{Итерационный метод решения линейных 


задач оптимального\newline управления с терминальными ограничениями}{3}


\tl{Васильева О.О., Васильев О.В.}


{К поиску равновесных управлений в дифференциальной\newline игре


$m$ лиц}{9}


\tl{Вознюк И.П.}


{Аппроксимационный алгоритм для задачи размещения на максимум\newline


с ограниченными объемами производств и поставок}{15}


\tl{Гимади Э.Х., Кайран Н.М., Сердюков А.И.}


{О разрешимости многоиндексной\newline аксиальной задачи о назначениях


на одноциклических подстановках}{21}


\tl{Емеличев В.А., Степанишина Ю.В.}


{Квазиустойчивость векторной нелинейной\newline траекторной задачи


с паретовским принципом оптимальности}{27}


\tl{Заботин И.Я.}


{Об устойчивости алгоритмов безусловной


минимизации псевдовыпуклых\newline функций}{33}


\tl{Заботин Я.И., Фукин И.А.}


{Об одной модификации метода сдвига штрафов для 


задач\newline нелинейного программирования}{49}


\tl{Коннов И.В.}


{Приближенные методы для прямо-двойственных 


вариационных неравенств\newline смешанного типа}{55}


\tl{Рязанцева И.П.}


{Метод итеративной регуляризации второго порядка для


выпуклых\newline задач условной минимизации}{67}


\tl{Срочко В.А., Пудалова Е.И.}


{Методы нелокального улучшения допустимых\newline управлений


в линейных задачах с запаздыванием}{78}








\bigskip





\centerline{КРАТКИЕ СООБЩЕНИЯ}


\bigskip





\tl{Кравцов М.К., Кравцов В.М., Лукшин Е.В.}


{О числе $r$-нецелочисленных


вершин\newline многогранника трехиндексной


аксиальной задачи выбора}{89}





\medskip\noindent{Содержание журнала за 2000 год\ \dotfill 93 \par





\vfill


\noindent\raisebox{2pt}{\copyright} Казанский госудаpственный унивеpситет


\end{document}


